\documentclass[letterpaper, 10 pt, conference]{ieeeconf}
\IEEEoverridecommandlockouts
\overrideIEEEmargins
\usepackage{amsmath,amssymb}
% ieeeconf (IEEEtran v1.6b) predefines \proof; let amsthm own it instead.
\let\proof\relax
\let\endproof\relax
\usepackage{amsthm}
\usepackage{cite}
\usepackage{graphicx}
\usepackage{booktabs}
\usepackage{orcidlink}
% orcidlink loads hyperref; keep links clickable but without boxes.
\hypersetup{hidelinks}
\usepackage[capitalise]{cleveref}

\newcommand{\R}{\mathbb{R}}
\newcommand{\norm}[1]{\left\lVert #1 \right\rVert}
\newcommand{\abs}[1]{\left\lvert #1 \right\rvert}

\newtheorem{theorem}{Theorem}
\newtheorem{lemma}{Lemma}
\newtheorem{corollary}{Corollary}
\newtheorem{assumption}{Assumption}
\newtheorem{remark}{Remark}
\crefname{assumption}{Assumption}{Assumptions}
\Crefname{assumption}{Assumption}{Assumptions}

\begin{document}
% Bibliography style controls defined in references.bib (IEEEexample:BSTcontrol).
\bstctlcite{IEEEexample:BSTcontrol}

\title{\LARGE \bf Finite-Time Centroid Tracking and Exponential Formation Stabilization in Arbitrary Swarm Configurations}

\author{Yash Purswani$^{*}$\orcidlink{0009-0002-9718-8160},
        Shivansh Agarwal$^{*}$\orcidlink{0009-0005-5756-2316}
        and Anuj Tiwari\orcidlink{0000-0003-3685-8888}%
\thanks{$^*$ Equal contribution and joint first authors. The authors are with the Department of Mechanical Engineering,
Indian Institute of Technology Madras, Chennai 600036, India.
{\tt\small me22b214@smail.iitm.ac.in},
{\tt\small me24s022@smail.iitm.ac.in},
{\tt\small anujt@iitm.ac.in}}%
}

\maketitle

\begin{abstract}
Certifying finite arrival times and convergence rates is essential for coordinating swarms in time-critical missions. However, these convergence properties generally depend on the interaction graph and feedback gains in existing methods, leaving exact acceptance deadlines uncertified for arbitrary time-varying formations. To address this limitation, this letter proposes a kinematic swarm model in which centered, time-varying stations are prescribed through individually assigned offsets with feedforward, driven by a shared finite-time centroid tracking feedback. It is shown that this shared feedback is dissipative for the relative formation error, which decouples the mission-level and agent-level channels in closed form. As a result, the swarm centroid is certified to track the reference in finite time, and individual station errors decay globally and exponentially from the initial condition. This decoupling also provides an explicit deadline certifying when every agent enters an arbitrary tolerance ball. Furthermore, the formation guarantee extends to limited sensing across switching proximity components. Numerical simulations confirm both certificates across rotating and reconfiguring formations.
\end{abstract}

\begin{keywords}
Control of networks, cooperative control, stability of nonlinear systems.
\end{keywords}

\section{Introduction}
\label{sec:intro}
Multi-robot teams routinely perform tasks difficult for individual platforms: aerial swarms map expansive terrain \cite{chung2018survey}, warehouse fleets sort inventory in parallel \cite{wurman2008coordinating}, and aerial manipulators jointly transport heavy loads \cite{michael2011cooperative}. Such a deployment is only useful once
the team behaves as one, following a common mission trajectory while
jointly holding a prescribed geometric arrangement -- the baseline of
a distributed sensor \cite{krieger2007tandem}, the pose of a jointly
carried load, or the corridor a team occupies through a confined
space. Such arrangements are frequently irregular rather than
symmetric, so a swarm model restricted to symmetric or graph-emergent
shapes cannot represent them.

In a growing share of these missions, the arrival is time-critical: a
docking maneuver has a window, a handoff or a grasp is safe only once
every agent is simultaneously within tolerance, and an intercept fails
outright if late. Such a guarantee must be computable before
deployment, from the initial condition alone, and must hold for each
agent individually in addition to the group average. Finite-time stability
\cite{haimo1986finite,bhat2000finite} certifies exact arrival with a time
bounded by the initial state, and has been incorporated into consensus and
formation designs \cite{wang2010finite,xiao2009finite}.

Formation control and trajectory tracking are each mature (\cref{sec:related}), but most schemes that combine them merge the reference and the formation shape into a single consensus-type error. Feasibility, convergence rate, and arrival time then all depend on the interaction graph, agent count, and gains; tracking and per-agent convergence are not certified independently, and existing finite-time results inherit the same dependence.

Fedele and D'Alfonso \cite{fedele2019kinematic} proposed a kinematic
swarm model in which a single fractional-power correction on the group
centroid renders its tracking finite-time, independently of the
interaction. The model has two limitations that matter here, i) the
desired formation is not independently prescribed, but emerges from the
pairwise interaction, so shapes that are not consistent equilibria of
the update law in \cite{fedele2019kinematic} cannot be realized; and, ii) only a swarm level convergence guarantee of the centroid is provided without individual-agent settling guarantee, since the agent formation error w.r.t. its desired position in the swarm is bounded only for
$t\ge\tau_c$, after the centroid has settled, leaving the preceding
transient uncertified.

This letter develops a kinematic swarm model that closes both gaps.
Each agent is assigned an independent, possibly time-varying offset
with feedforward about the moving centroid, subject only to the
centering condition that the offsets sum to zero, so any centered
formation -- symmetric or not -- is realizable and reconfigurable at
will. The corrector remains common to the swarm and acts only on the
centroid channel; it is shown to be dissipative for the formation
error as well, which decouples the group- and agent-level channels in
closed form and yields an individual-agent convergence certificate
holding from the initial instant. The contributions are:
\begin{enumerate}
\item a kinematic model with individually assigned, centered $C^1$
      offsets, so an arbitrary, asymmetric and time-varying formation
      is prescribed rather than emergent (\cref{sec:model});
\item finite-time centroid tracking with $\tau_c$ set by
      $\sigma(t_0)$ and $\beta$ alone, and not by $M$, $\nu$,
      $\varepsilon$, $n$ or $p_i(\cdot)$ (\cref{thm:centroid});
\item global exponential convergence to the stations from $t_0$, at a
      rate independent of $M$ (subject only to $M+M^\top\succeq0$), $\nu$, $\varepsilon$,
      $n$ and the formation geometry (\cref{thm:formation});
\item a closed-form acceptance deadline $T_{\mathrm{dock}}$ for any
      tolerance, computable at $t_0$ from $\norm{\delta(t_0)}$
      (\cref{cor:dock});
\item an extension to limited sensing, in which both formation results
      hold across any split of the proximity graph, which becomes
      complete after a computable time once the sensing radius exceeds
      the formation diameter (\cref{sec:connectivity}).
\end{enumerate}

\section{Background and Problem Statement}
\label{sec:related}

This section provides the relevant background on current state of formation control, finite time convergence literature, kinematic swarm models in the literature, which motivates the problem statement.

\noindent
\textbf{Formation control.} Classical formation strategies rely on position-, displacement-, distance-, or bearing-based feedback \cite{oh2015survey,su2026bearing}, virtual structures \cite{lewis1997high}, rigid frameworks \cite{anderson2008rigid}, or Laplacian maneuvers \cite{garciademarina2021distributed}. In each case, geometric shape and reference tracking are fundamentally coupled, leaving both the admissible formations and their convergence rates subordinate to the underlying interaction network.

\noindent
\textbf{Finite-, fixed-, and prescribed-time coordination.} Translating convergence rates into operational deadlines has driven substantial research. Finite-time consensus bounds depend directly on initial disagreement and algebraic connectivity \cite{wang2010finite}, while fixed-time methods \cite{polyakov2012nonlinear,zhou2022fixed,trinh2022finite} decouple the bound from initial states but retain scaling dependencies on network size, connectivity, and feedback gains. Prescribed-time controllers \cite{song2017time,song2023prescribed} assign exact settling times via time-varying gains, at the expense of terminal singularities on $[t_0, T)$ and elevated noise sensitivity. An explicit arrival certificate derived strictly from initial conditions via time-invariant feedback remains absent.

\noindent
\textbf{Swarm kinematic models.} One quantity does escape the
interaction. Building on the aggregation analysis of Gazi and Passino
\cite{gazi2003stability}, Fedele and D'Alfonso
\cite{fedele2019kinematic} drive the centroid with a common
finite-time corrector, alongside a distance-weighted interaction ruled
by a coupling matrix $M$, making the centroid's arrival
interaction-independent; extensions generalize $M$ to rotating, mixed,
discrete- and fixed-time variants
\cite{fedele2021mixing,fedele2020discrete}. The formation itself,
however, is still whatever the interaction produces, and the
agent-level guarantee is a hyper-ball of radius set by $M$, $\nu$,
$\varepsilon$ and $n$, entered only once the centroid has settled.

\noindent 
The existing approaches do not provide independent certificates for swarm-level tracking and agent-level formation keeping objectives simultaneously, which leads to the main research objectives for the current work as follows, 
\begin{enumerate}
\item \emph{Arbitrary formation prescription.} Specify an arbitrary,
      time-varying and asymmetric formation independently of the
      interaction graph and coupling matrix, rather than letting it
      emerge from the interaction;
\item \emph{Agent-level finite-time certification.} Certify, for every
      agent individually from $t_0$, an explicit deadline for entering
      an arbitrary tolerance of its station, independently of the
      interaction graph, the coupling gains and the swarm size.
\end{enumerate}

\section{Proposed Approach}

This section develops the kinematic model that addresses both
problems above: \cref{sec:model} defines the model together with the
assigned, centered stations; \cref{sec:trajectory} recovers the
finite-time centroid guarantee under the new model; and
\cref{sec:formation} establishes the corrector's dissipativity for the
formation error, from which the agent-level certificates follow;
\cref{sec:connectivity} relaxes the all-to-all sensing assumption.

\subsection{Kinematic Model}
\label{sec:model}

Consider $n\ge2$ agents in $\R^d$ with positions $x_i(t)$, $t\ge t_0$,
and centroid $c(t):=\frac1n\sum_{i=1}^n x_i(t)$; every agent knows the
position of every other, an assumption relaxed in \cref{sec:connectivity}.
The kinematic model of \cite{fedele2019kinematic} is extended by
assigning each agent a time-varying offset $p_i(t)$ together with its
feedforward $\dot p_i(t)$:
\begin{multline}
\dot x_i = -(x_i-p_i-r)+\dot r+\dot p_i-\zeta\\
-\frac{M}{n}\sum_{j=1}^n\underbrace{\frac{(x_i-p_i)-(x_j-p_j)}
{\norm{x_i-x_j}^\nu+\varepsilon}}_{=:\,\Gamma_{ij}},
\label{eq:model}
\end{multline}
where
\begin{itemize}
\item $r\in C^1$ is the reference for the centroid, and
      $\sigma:=c-r$ is the centroid error;
\item $p_i\in C^1$ places agent $i$'s \emph{assigned station} at
      $p_i(t)+r(t)$;
\item $\zeta_s=\mathrm{sign}(\sigma_s)\abs{\sigma_s}^\beta$,
      $s=1,\dots,d$, $\beta\in(0,1)$, is the finite-time corrector,
      common to all agents;
\item $M\in\R^{d\times d}$ with $M+M^\top\succeq0$ (i.e.\
      $v^\top Mv\ge0$ for all $v\in\R^d$; $M$ need not be symmetric),
      $\nu\in\mathbb{N}^+$ and $\varepsilon>0$ set the coupling gain,
      its decay with distance, and its regularization;
\item the coupling term $\Gamma_{ij}$ is skew-symmetric: its
      denominator $D_{ij}:=\norm{x_i-x_j}^\nu+\varepsilon$ is symmetric
      in $(i,j)$ and its numerator is skew-symmetric, so
      \begin{equation}
      \Gamma_{ij}=-\Gamma_{ji}.
      \label{eq:skewGamma}
      \end{equation}
\end{itemize}
For $p_i\equiv0$, Eq.~\eqref{eq:model} is exactly the model of
\cite[eq.~(1)]{fedele2019kinematic}. The offsets are constrained only
by the following assumption.
\begin{assumption}[Centering]
\label{ass:centering}
The offsets satisfy
\begin{equation}
\sum_{i=1}^n p_i(t)=0_d,\qquad\forall t\ge t_0.
\label{eq:centering}
\end{equation}
\end{assumption}
Since Eq.~\eqref{eq:centering} holds identically in $t$, it implies
$\sum_{i=1}^n\dot p_i(t)=0_d$ without further assumption, and it
places the mean station $\frac1n\sum_i(p_i+r)$ exactly on $r$.

\subsection{Trajectory Tracking}
\label{sec:trajectory}

\cref{thm:centroid} is based on \cite[Th.~1]{fedele2019kinematic}
with minor adaptations, \cref{ass:centering} removing the assigned
offsets from the centroid dynamics.

\begin{theorem}[Finite-Time Centroid Convergence \cite{fedele2019kinematic}]
\label{thm:centroid}
Under \cref{ass:centering}, the centroid error of Eq.~\eqref{eq:model}
satisfies $\sigma(t)\equiv0_d$ for all $t\ge\tau_c$, where
\begin{equation}
\tau_c = t_0+\max_{s=1,\dots,d}\frac{|\sigma_s(t_0)|^{1-\beta}}{1-\beta},
\label{eq:tauc}
\end{equation}
independent of $M,\nu,\varepsilon$ and of $p_i(\cdot)$.
\end{theorem}
\begin{proof}
Averaging Eq.~\eqref{eq:model} over $i=1,\dots,n$,
\[
\dot c(t) = -\frac1n\sum_i(x_i-p_i) + r + \dot r + \frac1n\sum_i\dot p_i
- \zeta - \frac{M}{n^2}\sum_i\sum_j\Gamma_{ij}.
\]
By the centering condition Eq.~\eqref{eq:centering},
$\frac1n\sum_i(x_i-p_i)=c(t)$ and $\frac1n\sum_i\dot p_i=0_d$; and by
the skew-symmetry Eq.~\eqref{eq:skewGamma},
$\sum_i\sum_j\Gamma_{ij}(t)=0_d$. Hence $\dot c=-c+r+\dot r-\zeta$,
and since $\sigma:=c-r$,
\begin{equation}
\dot\sigma(t) = -\sigma(t)-\zeta(t).
\label{eq:sigmadyn}
\end{equation}
Since $\zeta_s$ depends only on $\sigma_s$, Eq.~\eqref{eq:sigmadyn}
represents $d$ independent scalar equations, one per coordinate $s$,
\begin{equation}
\dot\sigma_s(t) = -\sigma_s(t)-\mathrm{sign}(\sigma_s(t))|\sigma_s(t)|^\beta,
\quad s=1,\dots,d.
\label{eq:sigmasdyn}
\end{equation}
For each $s$, consider the Lyapunov function
$V_s(t):=\frac12\sigma_s(t)^2$, the derivative of which along
Eq.~\eqref{eq:sigmasdyn} is exactly
\begin{align}
\dot V_s(t) &= \sigma_s\dot\sigma_s
= -\sigma_s(t)^2 - \sigma_s(t)\,\mathrm{sign}(\sigma_s(t))|\sigma_s(t)|^\beta
\notag\\
&= -\sigma_s(t)^2 - |\sigma_s(t)|^{1+\beta}.
\label{eq:Vsdot_exact}
\end{align}
Dropping the non-positive term $-\sigma_s(t)^2$ gives the bound
\begin{equation}
\dot V_s(t) \le -|\sigma_s(t)|^{1+\beta} = -2^\eta V_s(t)^\eta,
\label{eq:Vsdot}
\end{equation}
where the last equality rewrites $|\sigma_s|^{1+\beta}$ in terms of
$V_s$, using $|\sigma_s|=\sqrt{2V_s}$, and
$\eta:=\frac{1+\beta}{2}\in(\frac12,1)$. Finally, by the standard finite-time 
comparison argument \cite{bhat2000finite}, separating
variables and integrating \eqref{eq:Vsdot} from $t_0$ to $t$ gives
\begin{equation}
\begin{cases}
V_s^{1-\eta}(t) \le V_s^{1-\eta}(t_0) - 2^\eta(1-\eta)(t-t_0), & t\in[t_0,\tau_s]\\
V_s(t) = 0, & \forall t\ge\tau_s,
\end{cases}
\label{eq:Vsbound}
\end{equation}
where $t_0$ is the initial time and $\tau_s$ --- the instant at which
the bound in Eq.~\eqref{eq:Vsbound} reaches zero --- is
\begin{equation}
\tau_s = t_0 + \frac{V_s^{1-\eta}(t_0)}{2^\eta(1-\eta)}
= t_0 + \frac{|\sigma_s(t_0)|^{1-\beta}}{1-\beta}.
\label{eq:taus}
\end{equation}
Therefore, each component $\sigma_s(t)$ reaches zero in finite time,
no later than $\tau_s$. Setting $\tau_c:=\max_{s=1,\dots,d}\tau_s$
then gives $\sigma(t)\equiv0_d$ for all $t\ge\tau_c$, which is
Eq.~\eqref{eq:tauc}, independent of $M,\nu,\varepsilon,p_i(\cdot)$ since
these were eliminated already at Eq.~\eqref{eq:sigmadyn}.
\end{proof}
\subsection{Formation Tracking}
\label{sec:formation}

Define the formation error of agent $i$, $\delta_i:=x_i-p_i-r$, its
offset from its assigned station, stacked as
$\delta:=[\delta_1^\top,\dots,\delta_n^\top]^\top\in\R^{nd}$. Since
$(x_i-p_i)-(x_j-p_j)=\delta_i-\delta_j$, Eq.~\eqref{eq:model} gives
\begin{equation}
\dot\delta_i=-\delta_i-\zeta-\frac{M}{n}\sum_{j=1}^n\Gamma_{ij},\qquad
\Gamma_{ij}=\frac{\delta_i-\delta_j}{D_{ij}}.
\label{eq:error_dyn}
\end{equation}

\begin{lemma}[Centroid--Formation Identity]
\label{lem:identity}
Under \cref{ass:centering},
$\sum_{i=1}^n\delta_i(t)=n\,\sigma(t)$ for all $t\ge t_0$, and
consequently
\begin{equation}
-\sum_{i=1}^n\delta_i^\top(t)\zeta(t)=-n\sum_{s=1}^d|\sigma_s(t)|^{1+\beta}\le0.
\label{eq:dissipative}
\end{equation}
\end{lemma}
\begin{proof}
Summing $\delta_i=x_i-p_i-r$ over $i$ and using the centering
condition Eq.~\eqref{eq:centering},
\[
\sum_{i=1}^n\delta_i(t) = nc(t) - 0_d - nr(t) = n\sigma(t).
\]
Since $\zeta(t)$ carries no agent index, it factors out of the sum,
\begin{align*}
\sum_{i=1}^n\delta_i^\top(t)\zeta(t) &= n\sigma^\top(t)\zeta(t) \\
&= n\sum_{s=1}^d\sigma_s(t)\,\mathrm{sign}(\sigma_s(t))|\sigma_s(t)|^\beta \\
&= n\sum_{s=1}^d|\sigma_s(t)|^{1+\beta},
\end{align*}
which gives Eq.~\eqref{eq:dissipative}.
\end{proof}

\begin{theorem}[Global Exponential Formation Convergence]
\label{thm:formation}
Under \cref{ass:centering}, the solutions of Eq.~\eqref{eq:model} satisfy
$\norm{\delta(t)}\le\norm{\delta(t_0)}e^{-(t-t_0)}$ for all $t\ge t_0$,
independent of $M$ (subject only to $M+M^\top\succeq0$), $\nu$, $\varepsilon$, $n$ and the formation geometry; consequently
$\norm{\delta_i(t)}\le\norm{\delta(t)}$ for every $i$, and
$\delta_i(t)\to0_d$.
\end{theorem}
\begin{proof}
Consider the Lyapunov function
$V_\delta(t):=\frac12\norm{\delta(t)}^2=\frac12\sum_i\norm{\delta_i(t)}^2$.
Differentiating,
\[
\dot V_\delta(t) = \sum_i\delta_i^\top(t)\dot\delta_i(t).
\]
Substituting Eq.~\eqref{eq:error_dyn} and distributing $\delta_i^\top$,
\begin{align}
\dot V_\delta(t) &= -\sum_i\norm{\delta_i(t)}^2 -\sum_i\delta_i^\top(t)\zeta(t)
\notag\\
&\quad -\frac1n\sum_i\sum_j\delta_i^\top(t)M\Gamma_{ij}(t).
\label{eq:Vdot}
\end{align}
By Lemma~\ref{lem:identity}, the second term equals
$-n\sum_s|\sigma_s(t)|^{1+\beta}\le0$. For the third, let
$\Omega(t):=\sum_i\sum_j\delta_i^\top(t)M\Gamma_{ij}(t)$; relabeling
$i\leftrightarrow j$ and using $\Gamma_{ji}=-\Gamma_{ij}$ gives
$\Omega(t)=-\sum_i\sum_j\delta_j^\top(t)M\Gamma_{ij}(t)$, and averaging
the two expressions,
\begin{align}
\Omega(t) &= \frac12\sum_i\sum_j\big(\delta_i(t)-\delta_j(t)\big)^\top M\Gamma_{ij}(t)
\notag\\
&= \frac12\sum_i\sum_j\frac{\big(\delta_i(t)-\delta_j(t)\big)^\top M\big(\delta_i(t)-\delta_j(t)\big)}{D_{ij}(t)}.
\label{eq:Omega}
\end{align}
Since $M+M^\top\succeq0$, $(\delta_i-\delta_j)^\top M(\delta_i-\delta_j)\ge0$
(only the symmetric part of $M$ enters this quadratic form), and
$D_{ij}\ge\varepsilon>0$, so $\Omega(t)\ge0$ and the third term of
Eq.~\eqref{eq:Vdot} is non-positive. Discarding both non-positive terms,
\begin{equation}
\dot V_\delta(t) \le -\sum_i\norm{\delta_i(t)}^2 = -2V_\delta(t).
\label{eq:Vdotbound}
\end{equation}
Multiplying Eq.~\eqref{eq:Vdotbound} by $e^{2t}$ and using
$\frac{d}{dt}[e^{2t}V_\delta(t)] = e^{2t}\big(\dot V_\delta(t)+2V_\delta(t)\big)$,
\[
\frac{d}{dt}\big[e^{2t}V_\delta(t)\big] \le 0,
\]
so $e^{2t}V_\delta(t)$ is non-increasing; integrating from $t_0$ to
$t$ gives $V_\delta(t)\le V_\delta(t_0)e^{-2(t-t_0)}$, and since
$V_\delta=\frac12\norm{\delta}^2$,
\begin{equation}
\norm{\delta(t)}\le\norm{\delta(t_0)}e^{-(t-t_0)}, \qquad \forall t\ge t_0.
\label{eq:expbound}
\end{equation}
Since $\norm{\delta_i(t)}^2\le\sum_k\norm{\delta_k(t)}^2=\norm{\delta(t)}^2$
for every $i$, Eq.~\eqref{eq:expbound} gives
$\norm{\delta_i(t)}\le\norm{\delta(t)}\le\norm{\delta(t_0)}e^{-(t-t_0)}$,
and letting $t\to\infty$ gives $\delta_i(t)\to0_d$.
\end{proof}

\begin{remark}
Compared to \cref{thm:centroid}, note the difference in
convergence type: $\sigma$ reaches exactly zero at the finite time
$\tau_c$, since $\eta<1$ in Eq.~\eqref{eq:Vsdot} keeps the correction
strong all the way to the origin; here, once $\zeta$'s contribution
is absorbed by Lemma~\ref{lem:identity}, $\dot V_\delta\le-2V_\delta$
is a linear law, so $\delta(t)$ only approaches zero asymptotically.
\end{remark}

The bound holds from $t_0$: by \cref{lem:identity} the corrector is
dissipative for $\delta$ before the centroid settles, so neither
Eq.~\eqref{eq:expbound} nor the deadline below waits for $\tau_c$.

\begin{corollary}[Practical Finite-Time Convergence]
\label{cor:dock}
For any tolerance $\varepsilon_{\mathrm{tol}}>0$ (e.g.\ a physical
docking or gripper clearance margin),
$\norm{\delta_i(t)}\le\varepsilon_{\mathrm{tol}}$ for all $i$ and all
$t\ge T_{\mathrm{dock}}$, where
\begin{equation}
T_{\mathrm{dock}}=t_0+\max\!\Big(0,\ \ln\frac{\norm{\delta(t_0)}}{\varepsilon_{\mathrm{tol}}}\Big).
\label{eq:Tdock}
\end{equation}
\end{corollary}
\begin{proof}
By \cref{thm:formation},
\[
\norm{\delta_i(t)} \le \norm{\delta(t)} \le \norm{\delta(t_0)}e^{-(t-t_0)}
\]
for every agent $i$. Forcing the right-hand side below
$\varepsilon_{\mathrm{tol}}$ requires
\[
\norm{\delta(t_0)}e^{-(t-t_0)} \le \varepsilon_{\mathrm{tol}}
\quad\Longleftrightarrow\quad
t-t_0 \ge \ln\frac{\norm{\delta(t_0)}}{\varepsilon_{\mathrm{tol}}}
\]
whenever $\norm{\delta(t_0)}>\varepsilon_{\mathrm{tol}}$, and no wait
is needed if $\norm{\delta(t_0)}\le\varepsilon_{\mathrm{tol}}$; both
cases together give Eq.~\eqref{eq:Tdock}.
\end{proof}

\subsection{Extension to Limited Connectivity}
\label{sec:connectivity}

Following \cite[Sec.~V-A]{fedele2019kinematic}, suppose each agent senses only within radius $R_s>0$ and shares data within its connected component $\mathcal G_{q(i)}$ of size $n_q = \abs{\mathcal G_q}$. With component centroid $c_q = \frac{1}{n_q}\sum_{i\in\mathcal G_q}x_i$ and mean station offset $\bar p_q = \frac{1}{n_q}\sum_{i\in\mathcal G_q}p_i$, the component tracking error is $\tilde\sigma_q = c_q - \bar p_q - r$. Each agent restricts its corrector and coupling to its component:
\begin{equation}
\dot x_i = -(x_i-p_i-r) + \dot r + \dot p_i - \zeta(\tilde\sigma_{q(i)}) - \frac{M}{n_{q(i)}}\!\sum_{j\in\mathcal G_{q(i)}}\!\Gamma_{ij}.
\label{eq:model_conn}
\end{equation}

Both certificates survive this restriction because \cref{lem:identity} applies component-wise: summing $\delta_i$ over $\mathcal G_q$ yields $\sum_{i\in\mathcal G_q}\delta_i = n_q\tilde\sigma_q$ for any partition without requiring \cref{ass:centering}. Provided the topology switches finitely often on bounded intervals, the following guarantees hold:

\textbf{Formation invariance:} Within each component, the corrector remains dissipative. Because $\delta(t)$ is continuous across component mergers and splits, Eq.~\eqref{eq:expbound} and deadline $T_{\mathrm{dock}}$ \eqref{eq:Tdock} hold unconditionally for all $t\ge t_0$. In the decoupled limit $\mathcal G_q = \{i\}$, \eqref{eq:model_conn} simply reduces to $n$ independent finite-time station trackers.

\textbf{Centroid tracking:} On any dwell interval $[t_1, t_2)$ where components remain static, $\dot{\tilde\sigma}_q = -\tilde\sigma_q - \zeta(\tilde\sigma_q)$, certifying that each component error $\tilde\sigma_q$ vanishes by $t_1 + \max_s\abs{\tilde\sigma_{q,s}(t_1)}^{1-\beta}/(1-\beta)$. Under \cref{ass:centering}, $\sigma = \frac{1}{n}\sum_q n_q \tilde\sigma_q$, ensuring the global centroid reaches the reference once all components settle.

\textbf{Connectivity recovery:} Let $D_p := \sup_{t\ge t_0}\max_{i,j}\norm{p_i(t)-p_j(t)}$. Since $\norm{x_i - x_j} \le \sqrt2\norm{\delta} + D_p$, choosing $R_s > D_p$ guarantees the proximity graph becomes complete and remains connected for all $t \ge T_{\mathrm{conn}}$, where
\begin{equation}
T_{\mathrm{conn}} = t_0 + \max\!\Big(0,\ \ln\frac{\sqrt2\norm{\delta(t_0)}}{R_s - D_p}\Big).
\label{eq:Tconn}
\end{equation}
Beyond $T_{\mathrm{conn}}$, \cref{thm:centroid} applies globally. Network connectivity therefore accelerates transient coordination through coupling and ensures a unified centroid, but is not required to certify agent-level station acceptance.

\begin{figure*}[!t]
\centering
\includegraphics[width=\textwidth]{figures/f1_centroid.pdf}
\caption{Performance of the centroid channel across sliding exponents $\beta \in \{0.2, 0.5, 0.8\}$ alongside the linear baseline $\beta = 1$. (a) Convergence of the centroid tracking error $\Vert{}\sigma(t)\Vert{}$ plotted against the analytical settling-time bounds $\tau_c$ Eq.~\eqref{eq:tauc} (dashed lines). (b) Averaged control effort $\Vert{}u_{\text{avg}}(t)\Vert{} = \Vert{}\dot{c}(t)\Vert{}$ decaying to the feedforward baseline $\Vert{}\dot{r}\Vert{}$.}
\label{fig:centroid}
\end{figure*}

\section{Numerical Validation}
\label{sec:numerics}

To validate \cref{thm:centroid,thm:formation}
and Corollary~\ref{cor:dock}, the kinematic model Eq.~\eqref{eq:model} is
integrated numerically and the resulting trajectory is compared
against the analytical bounds Eq.~\eqref{eq:tauc}, Eq.~\eqref{eq:expbound},
and Eq.~\eqref{eq:Tdock}.

All experiments use the same swarm of $n=10$ agents in $d=2$
dimensions, started from the same initial positions and tracking
$r(t)=[t,\,0]^\top$ with the parameters of \cref{tab:params}. Three
experiments are reported: the centroid channel for several exponents
$\beta$ (\cref{sec:num_centroid}), the formation channel on a ring
template under several couplings (\cref{sec:num_formation}), and a
reconfiguration through four templates (\cref{sec:num_morph}).

% Solver details commented out for space; the same information is in
% Table I (Integrator, Arrival event).
% Integration uses LSODA at tight tolerances, since at loose ones the
% non-smooth term $\zeta$ leaves an error floor that can be mistaken for
% a genuine bound. A component $\sigma_s$ is declared to have arrived
% once $\abs{\sigma_s}=10^{-8}$ and is held at zero thereafter, which is
% exact because $\sigma_s=0$ is invariant.

\begin{table}[t]
\centering
\caption{Simulation setup.}
\label{tab:params}
\begin{tabular}{@{}ll@{}}
\toprule
Parameter & Value \\
\midrule
Agents, dimension ($n,d$) & $10,\ 2$ \\
Formation template & ring, radius $5$ \\
Rotation rate $\omega$ (time-varying run) & $0.5$~rad/s \\
Reconfiguration (Fig.~\ref{fig:morph}) & ring, wedge, star, arrow \\
Blend duration; start times & $4$~s;\ $t=14,26,38$~s \\
Reference $r(t)$ & $[t,\,0]^\top$ \\
Interaction matrix $M$ & $5.0\,I_2$ \\
Decay exponent $\nu$ & $2$ \\
Regularizer $\varepsilon$ & $1$ \\
Sliding exponent $\beta$ & $0.5$ (swept in Fig.~\ref{fig:centroid}) \\
Initial positions & $\mathcal N(0,25I_2)$, seed $42$ \\
Integrator & LSODA, rtol $10^{-10}$, atol $10^{-12}$ \\
Arrival event & $\abs{\sigma_s}=10^{-8}$, then latched \\
$\norm{\sigma(t_0)}$ & $1.5762$ \\
$\norm{\delta(t_0)}$ & $25.0657$ \\
\bottomrule
\end{tabular}
\end{table}

\begin{figure*}[!t]
\centering
\includegraphics[width=\textwidth]{figures/f2_morphing.pdf}
\caption{Trajectory profiles for a swarm of $n=10$ agents tracking a linear reference $r(t) = [t, 0]^\top$ while reconfiguring across four distinct centered templates: ring, wedge, star, and arrow. Dashed curves represent individual agent paths; hollow circles mark the time-varying target stations $p_i(t) + r(t)$ over template outlines, and filled dots indicate instant agent positions. Minimum-jerk blends over $4\text{ s}$ intervals preserve the centering condition $\sum p_i = 0$ throughout transitions, with agent-to-station slot assignment optimized via the Hungarian algorithm. A video of this simulation is available as {\color{blue}\href{https://drive.google.com/file/d/1XdVQPrmjcOjGwjsQR5pLfZK_i_-vPsrf/view?usp=sharing}{drive link}.}}
\label{fig:morph}
\end{figure*}


\subsection{Centroid Convergence and Its Control Cost}
\label{sec:num_centroid}

Fig.~\ref{fig:centroid}(a) shows $\norm{\sigma(t)}$ for
$\beta\in\{0.2,0.5,0.8\}$ and the linear baseline $\beta=1$. Each
fractional exponent reaches zero before its predicted $\tau_c$
(\cref{tab:beta}), larger $\beta$ arriving later, while the linear law
only decays exponentially.

\begin{table}[t]
\centering
\caption{Predicted settling time Eq.~\eqref{eq:tauc} and averaged control
effort for each exponent.}
\label{tab:beta}
\begin{tabular}{@{}cccc@{}}
\toprule
$\beta$ & $\tau_c$ (s) & peak $\norm{u_{\mathrm{avg}}}$
& $\int\norm{u_{\mathrm{avg}}-\dot r}\,dt$ \\
\midrule
$0.2$ & $1.51$ & $2.49$ & $1.594$ \\
$0.5$ & $2.25$ & $2.56$ & $1.580$ \\
$0.8$ & $5.24$ & $2.63$ & $1.576$ \\
$1$ & $\infty$ & $2.68$ & $1.576$ \\
\bottomrule
\end{tabular}
\end{table}

Fig.~\ref{fig:centroid}(b) shows the averaged command
$\norm{u_{\mathrm{avg}}}=\norm{\dot c}$. Averaging Eq.~\eqref{eq:model}
gives $\dot c=-\sigma+\dot r-\zeta(\sigma)$, since the coupling cancels
by Eq.~\eqref{eq:skewGamma} and the $\dot p_i$ by Eq.~\eqref{eq:centering}, so
the effort depends on the centroid error alone: it is largest at
$t_0$, decays as $\sigma$ closes, and settles on the feedforward level
$\norm{\dot r}$ once the centroid arrives. It dips below that level
while the corrector opposes the reference velocity, here because the
centroid starts ahead of $r$, and the jumps mark the arrival of
individual components $\sigma_s$. Since
$\abs{\sigma_s}^\beta<\abs{\sigma_s}$ for $\abs{\sigma_s}>1$ and the
reverse for $\abs{\sigma_s}<1$, a smaller $\beta$ commands less while an error
component exceeds one and more once it falls below, which is why the
fractional laws start lower yet arrive first.. The corrective effort
$\int\norm{u_{\mathrm{avg}}-\dot r}\,dt$ is nearly the same for every
$\beta$ (\cref{tab:beta}): it is the length of the path of $\sigma$ to the
origin, essentially the distance $\norm{\sigma(t_0)}$ whichever law
moves it.

\subsection{Formation Convergence and Acceptance Time}
\label{sec:num_formation}

\cref{tab:formation} shows the results on a  ring formation template using \cref{tab:params} config under several couplings. They include a
non-symmetric $M=5I_2+5J$, with $J$ the rotation by $90^\circ$, and the
ring rotating rigidly at $\omega=0.5$~rad/s, which exercises the
$\dot p_i$ feedforward and preserves \cref{ass:centering} since, with
$R(t)$ the rotation by $\omega t$, $\sum_iR(t)p_i=R(t)\sum_ip_i=0_d$. For each run, $\hat\lambda$ is the
largest $\lambda$ with
$\norm{\delta(t)}\le\norm{\delta(t_0)}e^{-\lambda(t-t_0)}$ while
$\norm{\delta}\ge10^{-6}$, which \cref{thm:formation} bounds below by
one, and the last two columns give the time after which
$\max_i\norm{\delta_i}\le\varepsilon_{\mathrm{tol}}$, against
$T_{\mathrm{dock}}$ of \cref{cor:dock}, common to all rows because
$\norm{\delta(t_0)}$ is. Every run meets both bounds. Stronger coupling
speeds convergence without
entering the certificate, and the uncoupled case attains the unit
rate, so the bound is tight when the interaction is off. Both
deadlines hold from $t_0$ rather than from $\tau_c$, as
\cref{lem:identity} supplies.

\begin{table}[t]
\centering
\caption{Formation convergence and acceptance time for several
couplings.}
\label{tab:formation}
\begin{tabular}{@{}lccc@{}}
\toprule
$M$ & $\hat\lambda$ & $\varepsilon_{\mathrm{tol}}=0.1$ (s)
& $\varepsilon_{\mathrm{tol}}=0.01$ (s) \\
\midrule
Bound Eq.~\eqref{eq:expbound}, Eq.~\eqref{eq:Tdock} & $1$ & $5.52$ & $7.83$ \\
\midrule
$0$ & $1.00$ & $5.06$ & $7.37$ \\
$5I_2$ & $1.14$ & $4.06$ & $6.08$ \\
$5I_2$, rotating & $1.13$ & $4.34$ & $6.36$ \\
$25I_2$ & $1.56$ & $2.55$ & $3.98$ \\
$5I_2+5J$ & $1.13$ & $4.08$ & $6.10$ \\
\bottomrule
\end{tabular}
\end{table}

\subsection{Formation Reconfiguration}
\label{sec:num_morph}

Fig.~\ref{fig:morph} illustrates the swarm transitioning across four distinct configurations: a ring, a filled wedge, a five-pointed star, and an arrow aligned with $r$. These geometries encompassing convex, non-convex, and filled lattice templates—are admissible because \cref{ass:centering} imposes no structural constraints beyond centering. Consecutive templates are interpolated via minimum-jerk blends $p_i(t)=(1-b(t))p_i^k + b(t)p_i^{k+1}$ with $b(t)\in[0,1]$, preserving Eq.~\eqref{eq:centering} identically and providing $\dot p_i$ in closed form. Because certificates \eqref{eq:tauc} and \eqref{eq:expbound} are independent of $p_i(\cdot)$, neither convergence guarantee is affected during reconfiguration.

Station allocation is resolved via the Hungarian algorithm \cite{kuhn1955hungarian} by minimizing $\sum_i\lVert x_i(t_0)-r(t_0)-p^{0}_{\pi(i)}\rVert^2$ at $t_0$ and $\sum_i\lVert p^{k}_i-p^{k+1}_{\pi(i)}\rVert^2$ at each template change. Station paths move synchronously along straight lines relative to $r$, guaranteeing collision-free virtual stations under sufficient endpoint clearance \cite{turpin2014capt}. In simulation, this assignment yields a minimum inter-agent clearance of $1.50$, compared to $0.006$ under naive index ordering. While effective in practice, this remains an allocation heuristic; actual agents follow stations within transient bounds $\delta_i$, so separated virtual targets do not strictly prevent physical collisions.

\section{Conclusion}
This letter shows that kinematic swarms can achieve simultaneous trajectory tracking and formation keeping with independent convergence guarantees. The core mechanism is a shared finite-time centroid feedback that dissipates collective formation errors. This provides two network-independent certificates -- i) the swarm centroid converges to the reference trajectory in finite time, and ii) individual station errors decay exponentially from the initial instant to meet an explicit acceptance deadline, which are confirmed by numerical simulations.  Independent of the network structure, coupling parameters can be tuned strictly for inter-agent spacing without invalidating the arrival time bounds. Furthermore, these guarantees hold under a limited communication radius across switching proximity components. %Numerical simulations confirmed both bounds, verifying that the swarm rotates and reconfigures through arbitrary geometries while respecting predicted deadlines. 
However, these results have certain limitations such as, practical implementation relies on distributed centroid estimation (e.g. dynamic average consensus), with local damping setting the certified exponential rate, and disturbances and actuation limits need to be incorporated.

\noindent
\textbf{Declaration on Generative AI Use:} The authors used generative AI tools, including ChatGPT and Claude, solely for the purpose of grammar checking and language refinement. All technical content, proofs, analysis, results, and conclusions were developed and verified by the authors.

\bibliographystyle{IEEEtran}
\bibliography{references}

\end{document}